% =====================================================================
\subsection{ETSI GS QKD 015: il controller SDN e la rete SD-QKD}
\label{sec:qkd015}
% =====================================================================

La GS~QKD~015 (V1.1.1 nel 2021, V2.1.1 nell'aprile 2022)~\cite{qkd015}
definisce l'interfaccia di controllo tra un controller SDN e i nodi di una
rete QKD, cioè la \emph{southbound interface} del piano di controllo QKD.
La seconda edizione, quella considerata qui, è stata allineata alla
GS~QKD~018: i moduli YANG della 018 importano direttamente quelli della 015.
Capire bene la 015 è quindi il prerequisito per capire la 018.

La sezione procede in quattro passi: perché serve un controllo software;
quali sono, una per una, le componenti di una rete SD-QKD; come è fatto in
generale un modello YANG; come la 015 usa YANG per descrivere il nodo, con
i workflow e le considerazioni di sicurezza che ne derivano.

% ---------------------------------------------------------------------
\subsubsection{Perché un controllo software-defined}

La specifica parte da una constatazione: integrare la QKD in una rete di
telecomunicazioni è oggi un'operazione manuale. Ogni nuovo sistema QKD
viene installato a mano, di solito su una fibra scura dedicata che collega
due siti, formando un link fisso. Questo approccio ha tre limiti. Non
scala, perché ogni modifica richiede intervento umano. Spreca risorse,
perché la capacità di chiave di un link non può essere spostata dove serve.
Infine non dialoga con la rete ottica, che invece è sempre più spesso
disaggregata e controllata da software.

L'architettura proposta permette a un controller di orchestrare
centralmente le risorse QKD: ottimizzare l'allocazione di chiave per link
in base alla domanda, creare automaticamente collegamenti diretti o
multi-hop, coordinare canali quantistici e classici anche quando
condividono la stessa fibra. In particolare l'astrazione definita dalla
specifica mira a:
\begin{itemize}
  \item gestire centralmente le risorse QKD in base a domande, capacità e
        stato della rete;
  \item aggregare i sistemi QKD di uno stesso perimetro sotto un'unica
        gestione delle chiavi, per riconoscere meglio le domande;
  \item nascondere la complessità dei parametri di basso livello di ogni
        tecnologia, delegandoli al nodo;
  \item ottimizzare la distribuzione dei link in base alla QoS delle
        applicazioni e alle statistiche di generazione;
  \item coordinare canali quantistici e classici, condividano o meno lo
        stesso mezzo fisico.
\end{itemize}
Sono, in forma tecnica, le esigenze emerse dai casi UC2, UC3 e UC4 della
GS~QKD~002.

% ---------------------------------------------------------------------
\subsubsection{Le componenti di una rete SD-QKD}

Conviene costruire la rete dal basso, aggiungendo un'astrazione alla volta.
La Fig.~\ref{fig:sdqkd} mostra il risultato.

\paragraph{1. Modulo, sistema e link QKD.}
L'unità fisica di base è il \emph{modulo QKD}: l'insieme di hardware e
software che esegue una parte di un protocollo QKD. Due moduli collegati
formano un \emph{sistema QKD}, e il collegamento tra loro (canale
quantistico più canale classico) è il \emph{QKD link}. Fin qui siamo nel
quantum layer della Fig.~\ref{fig:itu}: i moduli sono specifici del vendor
e della tecnologia (variabili discrete o continue, codifica, lunghezza
d'onda), e non sono interoperabili tra loro.

\paragraph{2. Interfaccia QKD.}
Il controller non ha bisogno di sapere come funziona internamente un
modulo, ma solo che cosa offre alla rete. La 015 introduce quindi
un'astrazione.

\begin{definition}[Interfaccia QKD]
Un'\emph{interfaccia QKD} (\texttt{qkdi}) è la rappresentazione astratta di
un modulo QKD all'interno di un nodo: ne espone solo gli attributi rilevanti
per la rete, come tecnologia, modello, intervallo di lunghezze d'onda,
attenuazione massima supportata e punto di attacco a uno switch ottico.
\end{definition}

L'analogia è quella delle porte di uno switch: un nodo con tre moduli QKD
appare al controller come un apparato con tre interfacce. Poiché moduli di
vendor diversi non interoperano, il modello impone di dichiarare tecnologia
e modello di ciascuna interfaccia: il controller abbinerà solo interfacce
compatibili.

\paragraph{3. Nodo SD-QKD.}
Le interfacce che stanno nello stesso luogo sicuro vengono raggruppate.

\begin{definition}[Nodo SD-QKD]
Un \emph{nodo SD-QKD} è l'aggregazione di uno o più moduli QKD situati in
un luogo sicuro e gestiti da un unico \emph{SDN agent}, che li espone al
controller come un singolo elemento di rete con più interfacce.
\end{definition}

La specifica fissa una serie di requisiti che danno al nodo la sua forma.
Un nodo SD-QKD \emph{deve}: comprendere almeno un modulo QKD, esposto come
interfaccia; trovarsi in un luogo sicuro (un luogo può ospitare più nodi);
contenere un \emph{sistema di gestione delle chiavi} che aggrega il
materiale prodotto da tutte le associazioni di chiave; offrire alle
applicazioni un \emph{unico punto di accesso} alle chiavi tramite API;
collegarsi ad almeno un controller SDN con protocolli standard, per
discovery, controllo e telemetria. Inoltre \emph{dovrebbe} esporre al
controller le informazioni su applicazioni, interfacce, link e requisiti
del canale classico (per esempio l'attenuazione accettabile).

All'interno del nodo si riconoscono quindi quattro componenti logiche:
\begin{itemize}
  \item i \textbf{moduli QKD}, che generano chiave con i nodi remoti;
  \item il \textbf{key manager} (KMS), che riceve le chiavi dai moduli,
        le memorizza in uno o più \emph{key store} logici, le inoltra in
        caso di relay e registra le applicazioni con la loro QoS;
  \item l'\textbf{interfaccia verso le applicazioni}, cioè la key
        delivery API (GS~QKD~004 o~014), raggiungibile a un indirizzo
        IP e porta noti;
  \item l'\textbf{SDN agent}, che traduce lo stato di tutto ciò in un
        modello YANG, lo espone al controller, riceve configurazioni e
        invia notifiche.
\end{itemize}

% Figura: rete SD-QKD e struttura interna di un nodo (adattata da GS QKD 015, Fig. 1)
\begin{figure}[t]
\centering
\begin{tikzpicture}[
  comp/.style={draw,rounded corners,align=center,font=\scriptsize,
               minimum height=6mm},
  small/.style={draw,circle,fill=green!12,minimum size=9mm,font=\scriptsize,
                align=center}]
  % controller
  \node[box,fill=blue!8,minimum width=4cm] (ctl) at (-0.6,4.3) {SDN controller (SD-QNC)};
  % nodo dettagliato
  \node[comp,fill=yellow!20,minimum width=1.4cm] (m1) at (-2.2,0) {modulo\\QKD 1};
  \node[comp,fill=yellow!20,minimum width=1.4cm] (m2) at (-0.6,0) {modulo\\QKD 2};
  \node[comp,fill=yellow!20,minimum width=1.4cm] (m3) at (1.0,0) {modulo\\QKD 3};
  \node[comp,fill=green!12,minimum width=4.6cm] (kms) at (-0.6,1.0) {key manager (KMS, key store)};
  \node[comp,fill=blue!6,minimum width=1.9cm] (agent) at (-1.95,2.0) {SDN agent};
  \node[comp,fill=gray!10,minimum width=2.4cm] (app) at (0.75,2.0) {applicazioni\\(QKD 004/014)};
  \begin{scope}[on background layer]
    \node[draw,dashed,rounded corners,fill=gray!3,inner sep=6pt,
          fit=(m1)(m3)(kms)(agent)(app),
          label={[font=\scriptsize\itshape]below:nodo SD-QKD (perimetro sicuro)}] (node) {};
  \end{scope}
  % nodi remoti
  \node[small] (r1) at (4.6,2.6) {nodo\\B};
  \node[small] (r2) at (4.4,-0.2) {nodo\\C};
  \node[small] (r3) at (-5.2,0.6) {nodo\\D};
  % collegamenti
  \draw[thick] (m3.east) -- node[lbl,above]{PHYS} (r2);
  \draw[thick] (m1.west) -- (r3);
  \draw[thick] (r1) -- (r2);
  \draw[thick,dotted,red!70!black] (kms.east) -- node[pos=0.55,font=\scriptsize,fill=white,inner sep=1pt]{VIRT via C} (r1);
  \draw[<->] (m1) -- (m1|-kms.south);
  \draw[<->] (m2) -- (m2|-kms.south);
  \draw[<->] (m3) -- (m3|-kms.south);
  \draw[<->] (agent.south) -- (agent.south|-kms.north);
  \draw[<->] (app.south) -- (app.south|-kms.north);
  \draw[<->,dashed,blue!60!black] (agent.north) -- node[pos=0.5,left,font=\scriptsize]{GS QKD 015} (ctl);
  \draw[<->,dashed,blue!60!black] (r1) -- (ctl);
  \draw[<->,dashed,blue!60!black] (r3) -- (ctl);
\end{tikzpicture}
\caption[Rete SD-QKD e struttura interna di un nodo]{Una rete SD-QKD: i nodi sono collegati da link fisici (linee
continue) e virtuali (punteggiata) e ciascuno è controllato dal controller
tramite il proprio SDN agent (tratteggio). Il nodo in primo piano mostra le
componenti interne: i moduli QKD, esposti come interfacce, il key manager,
l'interfaccia verso le applicazioni e l'agent. Adattata
da~\cite[Fig.~1]{qkd015}.}
\label{fig:sdqkd}
\end{figure}


\paragraph{4. Key association link.}
Tra due nodi la grandezza interessante non è il canale quantistico in sé
ma il fatto che i due nodi condividano chiave. La 015 modella proprio
questo.

\begin{definition}[Key association link]
Un \emph{key association link} (\texttt{qkd\_link}) è un'associazione
logica di chiave tra due nodi SD-QKD remoti. È di tipo \texttt{PHYS}
(diretto) se esiste un canale quantistico tra i due nodi, di tipo
\texttt{VIRT} (virtuale) se la chiave è ottenuta tramite key relay
attraverso nodi fidati intermedi.
\end{definition}

Un link fisico corrisponde a un QKD link della terminologia ITU-T; un link
virtuale corrisponde invece a una catena di KM link sopra più QKD link. Due
dettagli rendono il modello coerente. Il primo: un link virtuale
\emph{consuma} chiave dai link fisici su cui poggia, e per questo viene
registrato come \emph{applicazione interna} con una banda riservata
(\texttt{virt\_bandwidth}). Il secondo: ogni nodo intermedio di un link
virtuale conosce il nodo precedente e il successivo
(\texttt{virt\_prev\_hop}, \texttt{virt\_next\_hop}), quindi il cammino di
relay è configurato hop per hop dal controller. Perché un nodo possa fare
da relay deve dichiararlo con la capability
\texttt{key\_relay\_mode\_enable}.

\paragraph{5. Applicazioni.}
Dal punto di vista del nodo, un'applicazione è qualunque entità che chiede
chiavi al key manager.

\begin{definition}[Applicazione QKD]
Un'\emph{applicazione QKD} (\texttt{qkd\_app}) è un'istanza in esecuzione
che consuma chiavi da un nodo SD-QKD. È \emph{esterna}
(\texttt{CLIENT}) se si tratta di un utente, un HSM, un cifratore, una
funzione di rete virtuale o un protocollo di sicurezza; è \emph{interna}
(\texttt{INTERNAL}) se usa le chiavi per il funzionamento della rete
stessa, come l'autenticazione o il relay di un link virtuale.
\end{definition}

Ogni applicazione dichiara una qualità del servizio: banda massima e
minima di chiave (\texttt{max\_bandwidth}, \texttt{min\_bandwidth}, in
bit/s), \texttt{jitter} massimo della consegna in ms, tempo di vita
\texttt{ttl} di una chiave inutilizzata nel key store e la possibilità di
condividere cammini o chiavi tra più client. È la risposta tecnica
all'esigenza di QoS differenziata del caso UC2.

\paragraph{6. Controller SDN.}
Al vertice c'è il controller (nella 015 anche \emph{SD-QNC},
\emph{Software-Defined QKD Network Controller}), che ha la vista globale del
dominio e svolge tre ruoli. Fa da \emph{registro centrale} delle
applicazioni, così che un'applicazione debba conoscere solo il proprio key
manager e non la topologia della rete. Pianifica e crea i \emph{link},
fisici o virtuali, in base a domande e statistiche. Coordina infine il
proprio dominio con il controller della rete ottica (\emph{SD-ONC}) quando
serve un nuovo cammino ottico per un canale quantistico.

\paragraph{Che cosa non sale al controller.}
L'informazione che attraversa il piano di controllo è volutamente minima:
identificativi, stati, tassi, parametri fisici. Il materiale di chiave
\emph{non lascia mai il perimetro del nodo}. Per questo, osserva la
specifica, introdurre SDN in una rete a nodi fidati non aggiunge rischi
nuovi oltre a quelli già noti, come il denial of service.

% ---------------------------------------------------------------------
\subsubsection{Anatomia di un modello YANG}
\label{sec:yang-anatomia}

Sia la 015 sia la 018 sono scritte in YANG~\cite{rfc7950}, un linguaggio
nato in IETF per descrivere i dati di configurazione e di stato degli
apparati di rete e le operazioni che si possono invocare su di essi. Un
modello YANG non è codice eseguibile: è uno \emph{schema}, come un XSD per
XML o uno schema JSON, che stabilisce quali dati esistono, di che tipo sono,
come sono annidati e chi può modificarli. Su quello schema poi lavorano i
protocolli NETCONF e RESTCONF. Vediamo i costrutti principali, usando
esempi tratti dai file della 015~\cite{qkd015yang}.

\paragraph{Il modulo e la sua intestazione.}
L'unità di base è il \emph{modulo}. L'intestazione dichiara la versione
del linguaggio, un \texttt{namespace} globale univoco, un \texttt{prefix}
breve con cui gli altri moduli lo citano, i moduli importati e le revisioni.
\begin{lstlisting}[language=yang]
module etsi-qkd-sdn-node {
  yang-version "1";
  namespace "urn:etsi:qkd:yang:etsi-qkd-node";
  prefix "etsi-qkdn";
  import ietf-yang-types   { prefix "yang"; }
  import ietf-inet-types   { prefix "inet"; }
  import etsi-qkd-node-types { prefix "etsi-qkdn-types"; }
  organization "ETSI ISG QKD";
  revision "2022-01-30" { reference "ETSI GS QKD 015 V2.1.1"; }
  ...
}
\end{lstlisting}
La 015 è divisa in due moduli: \texttt{etsi-qkd-node-types} contiene solo
tipi, \texttt{etsi-qkd-sdn-node} contiene il modello vero e proprio. È una
scelta comune: separare i tipi permette ad altri moduli, come quelli della
018, di riusarli senza importare tutto il resto.

\paragraph{Tipi: built-in, typedef e identity.}
YANG offre tipi predefiniti (\texttt{boolean}, \texttt{string},
\texttt{uint32}, \texttt{decimal64}, \texttt{enumeration}\dots) e tipi
comuni definiti dall'IETF, come \texttt{yang:uuid} o
\texttt{yang:date-and-time}. Con \texttt{typedef} si definiscono tipi
derivati con un nome proprio. Per valori simbolici estendibili si usano le
\texttt{identity}: si dichiara un'identità base e poi identità che la
derivano; un campo di tipo \texttt{identityref} accetta qualunque identità
derivata dalla base. È così che la 015 definisce i tipi di link:
\begin{lstlisting}[language=yang]
identity QKD-LINK-TYPES { description "QKD key association link types."; }
identity VIRT { base QKD-LINK-TYPES; description "Virtual Link."; }
identity PHYS { base QKD-LINK-TYPES; description "Physical Link."; }

typedef qkd-link-types {
  type identityref { base QKD-LINK-TYPES; }
}
\end{lstlisting}
Il vantaggio rispetto a un'\texttt{enumeration} è che un altro modulo può
aggiungere, per esempio, un nuovo tipo di link senza modificare l'originale.
Allo stesso modo sono definiti le tecnologie (\texttt{CV-QKD},
\texttt{DV-QKD}\dots), gli stati di nodo (\texttt{NEW},
\texttt{OPERATING}, \texttt{DOWN}, \texttt{FAILURE}, \texttt{OUT}), i ruoli
(\texttt{TRANSMITTER}, \texttt{RECEIVER}, \texttt{TRANSCEIVER}) e i tipi
di applicazione (\texttt{CLIENT}, \texttt{INTERNAL}).

\paragraph{Nodi dati: container, list, leaf, leaf-list.}
I dati sono organizzati ad albero con quattro costrutti:
\begin{itemize}
  \item \texttt{leaf}: un singolo valore tipizzato (per esempio
        \texttt{qkdn\_id});
  \item \texttt{leaf-list}: una sequenza di valori dello stesso tipo
        (per esempio \texttt{virt\_next\_hop}, che ammette più successori);
  \item \texttt{container}: un nodo che raggruppa altri nodi, senza
        valore proprio (per esempio \texttt{qkdn\_capabilities});
  \item \texttt{list}: una collezione di elementi strutturati, ciascuno
        identificato da una \texttt{key} (per esempio la lista
        \texttt{qkd\_link} con chiave \texttt{qkdl\_id}).
\end{itemize}

\paragraph{Configurazione e stato.}
Ogni nodo è di \emph{configurazione} (\texttt{config true}, il default:
il client può scriverlo) oppure di \emph{stato operativo}
(\texttt{config false}: lo produce l'apparato e il client può solo
leggerlo). Nella 015, per esempio, \texttt{qkdl\_enable} è configurabile,
mentre \texttt{qkdn\_status}, \texttt{skr} ed \texttt{eskr} sono stato.
Questa distinzione riflette la divisione tra ciò che il controller
\emph{decide} e ciò che il nodo \emph{misura}.

\paragraph{Riuso ed estensione: grouping, uses, augment.}
Un \texttt{grouping} è un blocco di nodi riutilizzabile, che viene
copiato dove compare \texttt{uses}. Il modello della 015 è costruito quasi
tutto così: \texttt{app\_basic}, \texttt{app\_details}, \texttt{app\_qos},
\texttt{qkd\_links} e così via sono grouping poi assemblati nel container
\texttt{qkd\_node}. L'\texttt{augment} fa l'operazione opposta: aggiunge
nodi a un albero definito \emph{in un altro modulo}, senza toccarlo. È lo
strumento che la 018 userà per arricchire la propria topologia con i dati
di inventario.

\paragraph{Riferimenti e vincoli.}
Un \texttt{leafref} è un campo il cui valore deve coincidere con un valore
esistente altrove nell'albero, come una chiave esterna in un database. Le
istruzioni \texttt{when} e \texttt{must} aggiungono condizioni, per esempio
che un campo esista solo se il link è di un certo tipo. La 015 indica
proprio così che i campi \texttt{virt\_*} valgono solo per link virtuali e
i campi \texttt{phys\_*} solo per quelli fisici.

\paragraph{Operazioni e notifiche: rpc e notification.}
Oltre ai dati, un modulo può definire operazioni invocabili dal client
(\texttt{rpc}, con \texttt{input} e \texttt{output}) e messaggi che il
server invia spontaneamente (\texttt{notification}). La 015 definisce
soprattutto notifiche; la 018, come vedremo, definisce anche RPC.

\paragraph{Leggere un albero.}
Lo strumento \texttt{pyang} produce una vista ad albero di un modulo, che
useremo nel seguito. La notazione è compatta: \texttt{rw} indica un nodo di
configurazione, \texttt{ro} un nodo di stato, \texttt{?} un nodo
facoltativo, \texttt{*} una lista o leaf-list, la chiave di una lista
compare tra parentesi quadre e \texttt{-{}-{}-n} indica una notifica.

\paragraph{Dai dati al messaggio.}
Lo stesso schema produce sia XML sia JSON. In JSON~\cite{rfc7951} il primo
livello è qualificato con il nome del modulo e i valori di tipo
\texttt{identityref} riportano il modulo che definisce l'identità. Un link
fisico esposto da un nodo apparirebbe così:
\begin{lstlisting}
{ "etsi-qkd-sdn-node:qkd_node": {
    "qkdn_id": "0f3c...-a1",
    "qkd_links": { "qkd_link": [ {
        "qkdl_id": "7b2e...-c4",
        "qkdl_type": "etsi-qkd-node-types:PHYS",
        "qkdl_local":  { "qkdn_id": "0f3c...-a1", "qkdi_id": 1 },
        "qkdl_remote": { "qkdn_id": "9d41...-e7", "qkdi_id": 2 },
        "phys_wavelength": 1550,
        "qkdl_performance": { "skr": 20000, "eskr": 12000,
                              "expected_consumption": 8000 } } ] } } }
\end{lstlisting}
(I valori sono illustrativi.)

% ---------------------------------------------------------------------
\subsubsection{Il modello del nodo SD-QKD}

Con questi strumenti il modello della 015 si legge facilmente. Il container
radice \texttt{qkd\_node} ha quattro sezioni, una per ciascuna componente
descritta sopra, più le notifiche:
\begin{lstlisting}
module: etsi-qkd-sdn-node
  +--rw qkd_node
     +--rw qkdn_id              yang:uuid
     +--ro qkdn_status?         etsi-qkdn-types:qkdn-status-types
     +--rw qkdn_version?        string
     +--rw qkdn_location_id?    string
     +--rw qkdn_capabilities
     |  +--rw link_stats_support?          boolean
     |  +--rw application_stats_support?   boolean
     |  +--rw key_relay_mode_enable?       boolean
     +--rw qkd_applications
     |  +--rw qkd_app* [app_id]
     |     +--rw app_id, app_type, server_app_id, client_app_id*, ...
     |     +--rw app_qos            (max/min_bandwidth, jitter, ttl, ...)
     |     +--rw backing_qkdl_id*, local_qkdn_id, remote_qkdn_id
     +--rw qkd_interfaces
     |  +--rw qkd_interface* [qkdi_id]
     |     +--rw qkdi_id, qkdi_model, qkdi_type
     |     +--rw qkdi_att_point     (device, port)
     |     +--rw qkdi_capabilities  (role, wavelength_range, max_absorption)
     |     +--ro qkdi_status
     +--rw qkd_links
        +--rw qkd_link* [qkdl_id]
           +--rw qkdl_id, qkdl_status, qkdl_enable, qkdl_type
           +--rw qkdl_local / qkdl_remote   (qkdn_id, qkdi_id)
           +--rw qkdl_applications*
           +--rw virt_prev_hop, virt_next_hop*, virt_bandwidth   (VIRT)
           +--rw phys_channel_att, phys_wavelength, phys_qkd_role (PHYS)
           +--rw qkdl_performance   (expected_consumption, skr, eskr,
                                     phys_perf* [perf_type]: QBER, SNR)
\end{lstlisting}

\paragraph{Identità del nodo.}
\texttt{qkdn\_id} è un UUID univoco nella rete; \texttt{qkdn\_status},
\texttt{qkdn\_version} e \texttt{qkdn\_location\_id} servono a inventario e
diagnostica.

\paragraph{Capability.}
Le tre capability dicono al controller che cosa il nodo è disposto a fare:
esporre statistiche sui link (SKR, consumo, stato, QBER), esporre
statistiche sulle applicazioni e fare da relay. Un nodo con
\texttt{key\_relay\_mode\_enable = false} non verrà mai scelto come nodo
intermedio di un link virtuale.

\paragraph{Applicazioni.}
Ogni \texttt{qkd\_app} lega un'applicazione al servizio di chiave:
\texttt{server\_app\_id} identifica l'entità che ha avviato la richiesta,
\texttt{client\_app\_id} le entità autorizzate a ricevere la stessa chiave,
\texttt{backing\_qkdl\_id} il link che fornisce le chiavi,
\texttt{local\_qkdn\_id} e \texttt{remote\_qkdn\_id} i due nodi coinvolti.
Finché il nodo remoto non è noto, il nodo locale non può servire chiavi.
Le statistiche (\texttt{app\_statistics}) riportano i bit consumati per
intervallo di tempo.

\paragraph{Interfacce.}
Ogni \texttt{qkd\_interface} descrive un modulo: modello, tecnologia, stato,
punto di attacco a un dispositivo ottico (\texttt{qkdi\_att\_point}) e
capability fisiche. È l'informazione che serve al controller per abbinare
due interfacce compatibili e chiedere alla rete ottica un cammino tra le
rispettive porte.

\paragraph{Link e prestazioni.}
Ogni \texttt{qkd\_link} riporta gli estremi, il tipo, lo stato, le
applicazioni servite e i parametri specifici del tipo. Il blocco
\texttt{qkdl\_performance} contiene le tre grandezze che guidano tutta la
pianificazione:
\begin{itemize}
  \item \texttt{skr}: tasso di generazione di chiave segreta, in bit/s;
  \item \texttt{expected\_consumption}: somma delle bande di tutte le
        applicazioni sul link;
  \item \texttt{eskr}: tasso \emph{effettivo}, cioè disponibile dopo il
        consumo interno.
\end{itemize}
In forma compatta, se $\mathcal{A}_\ell$ è l'insieme delle applicazioni sul
link $\ell$, ciascuna con banda $b_a$,
\begin{equation}
  \text{expected\_consumption}_\ell = \sum_{a\in\mathcal{A}_\ell} b_a ,
  \qquad
  \mathrm{ESKR}_\ell \approx \mathrm{SKR}_\ell - \text{expected\_consumption}_\ell .
  \label{eq:eskr015}
\end{equation}
La seconda relazione è una lettura intuitiva dei campi, non una formula
dello standard. Per i link fisici si aggiungono metriche come QBER e SNR.

% ---------------------------------------------------------------------
\subsubsection{Notifiche}

Interrogare continuamente i nodi sarebbe inefficiente. La 015 prevede
quindi notifiche YANG dal nodo al controller, divise per oggetto:
\begin{itemize}
  \item \textbf{applicazioni}: \texttt{sdqkdn\_application\_new} (una nuova
        applicazione chiede chiavi), \texttt{sdqkdn\_application\_qos\_update},
        \texttt{sdqkdn\_application\_disconnected};
  \item \textbf{interfacce}: \texttt{sdqkdn\_interface\_new},
        \texttt{sdqkdn\_interface\_down} (con un campo \emph{reason}),
        \texttt{sdqkdn\_interface\_out};
  \item \textbf{link}: \texttt{sdqkdn\_link\_down},
        \texttt{sdqkdn\_link\_perf\_update} (variazione significativa delle
        prestazioni), \texttt{sdqkdn\_link\_overloaded} (il link non regge
        la domanda).
\end{itemize}
A queste si aggiunge una struttura generica di eventi e allarmi, pensata
per accogliere nuovi tipi di evento senza cambiare il modello. Inoltre i
protocolli basati su YANG permettono di sottoscrivere qualunque dato del
datastore, così che ogni variazione possa essere segnalata in tempo reale.

% ---------------------------------------------------------------------
\subsubsection{I workflow}

La clausola 5 descrive tre scenari fondamentali, che mostrano le componenti
all'opera.

\paragraph{Registrazione di un'applicazione.}
È il caso più istruttivo, perché spiega perché il controller fa da registro
centrale (Fig.~\ref{fig:seq015}). Un'applicazione conosce solo l'indirizzo
del proprio key manager; non conosce il nodo dell'applicazione con cui
vuole comunicare.
\begin{enumerate}
  \item L'applicazione nel dominio di sicurezza $A$ si collega al proprio
        nodo tramite la key delivery API, inviando il proprio ID, l'ID
        dell'applicazione remota e la QoS richiesta.
  \item Il nodo $A$ non sa dove si trovi l'applicazione remota e inoltra la
        richiesta al controller, aggiungendo il proprio ID.
  \item Se l'applicazione remota non è ancora registrata, il controller
        risponde con un semplice riscontro.
  \item Nel dominio $B$ l'applicazione remota si collega al proprio nodo, che
        a sua volta informa il controller.
  \item Il controller riconosce la coppia, genera un identificativo unico
        di servizio (\texttt{app\_id}) e invia a entrambi i nodi tutte le
        informazioni necessarie: ID di applicazioni e nodi, \texttt{app\_id}
        e QoS.
  \item I nodi comunicano l'\texttt{app\_id} alle applicazioni e riservano
        chiavi per esse; da qui in poi le applicazioni chiedono chiavi
        usando quell'identificativo.
\end{enumerate}

% Figura: registrazione di un'applicazione (adattata da GS QKD 015, Fig. 2)
\begin{figure}[t]
\centering
\begin{tikzpicture}[font=\scriptsize]
  \def\xA{0}\def\xB{3.0}\def\xC{6.6}\def\xD{10.2}\def\xE{13.2}
  \node[actor,minimum width=2cm] at (\xA,0) {App A};
  \node[actor,minimum width=2cm,fill=green!8] at (\xB,0) {Nodo A};
  \node[actor,minimum width=2cm,fill=blue!8] at (\xC,0) {Controller};
  \node[actor,minimum width=2cm,fill=green!8] at (\xD,0) {Nodo B};
  \node[actor,minimum width=2cm] at (\xE,0) {App B};
  \foreach \x in {\xA,\xB,\xC,\xD,\xE} \draw[gray] (\x,-0.4) -- (\x,-6.6);
  \draw[msg] (\xA,-0.9) -- node[above]{ID, ID peer, QoS} (\xB,-0.9);
  \draw[msg] (\xB,-1.5) -- node[above]{\texttt{sdqkdn\_application\_new}} (\xC,-1.5);
  \draw[ret] (\xC,-2.1) -- node[above]{ack (peer ignoto)} (\xB,-2.1);
  \draw[msg] (\xE,-2.9) -- node[above]{ID, ID peer, QoS} (\xD,-2.9);
  \draw[msg] (\xD,-3.5) -- node[above]{\texttt{sdqkdn\_application\_new}} (\xC,-3.5);
  \draw[msg] (\xC,-4.1) -- ++(0.35,0) -- ++(0,-0.35) -- (\xC,-4.45);
  \node[right] at (\xC+0.4,-4.27) {abbina e crea \texttt{app\_id}};
  \draw[msg] (\xC,-5.1) -- node[above]{config.\ (\texttt{app\_id}, QoS)} (\xB,-5.1);
  \draw[msg] (\xC,-5.1) -- node[above]{config.\ (\texttt{app\_id}, QoS)} (\xD,-5.1);
  \draw[ret] (\xB,-5.9) -- node[above]{\texttt{app\_id}} (\xA,-5.9);
  \draw[ret] (\xD,-5.9) -- node[above]{\texttt{app\_id}} (\xE,-5.9);
\end{tikzpicture}
\caption[Registrazione di un'applicazione in una rete SD-QKD]{Registrazione di un'applicazione in una rete SD-QKD. Le
applicazioni conoscono solo il proprio key manager; è il controller ad
abbinare le due richieste e a creare l'identificativo di servizio.
Adattata da~\cite[Fig.~2]{qkd015}.}
\label{fig:seq015}
\end{figure}


\paragraph{Creazione di un link fisico.}
Il controller riceve un innesco, da una domanda rilevata sui nodi o dalla
propria NBI. Cerca nei due siti due interfacce libere e compatibili,
chiede al controller della rete ottica (SD-ONC) un cammino ottico
(\emph{lightpath}) con i vincoli del canale quantistico e, ricevuta la
conferma, configura le due interfacce sui rispettivi nodi. Se i due
controller sono integrati in uno solo, l'operazione è interna. Questo
workflow è il cuore della coesistenza tra QKD e rete ottica chiesta dal
caso UC4.

\paragraph{Creazione di un link virtuale.}
Se non ci sono interfacce compatibili o non esiste un lightpath adatto, il
controller ripiega su un link virtuale. Calcola un cammino sulla rete QKD
esistente, per esempio su un grafo pesato costruito con disponibilità di
chiave, tassi di generazione e consumi; configura ogni nodo intermedio in
modalità relay indicando precedente e successivo; infine configura gli
estremi con la banda da fornire. La specifica non prescrive l'algoritmo.
Una scelta naturale, coerente con la~\eqref{eq:eskr015}, è pesare ciascun
link con l'inverso della chiave disponibile e scartare i link che non
reggono la banda $r$ richiesta:
\begin{equation}
  w_\ell = \frac{1}{\mathrm{ESKR}_\ell}, \qquad
  p^\star = \arg\min_{p\,:\,\min_{\ell\in p}\mathrm{ESKR}_\ell \ge r}
            \;\sum_{\ell\in p} w_\ell .
  \label{eq:path015}
\end{equation}
La~\eqref{eq:path015} è un esempio nostro, non parte dello standard. La
specifica avverte che, se gestiti male, i link virtuali possono congestionare
i link fisici sottostanti: è proprio per questo che serve un controller con
la vista globale di consumi e disponibilità.

% ---------------------------------------------------------------------
\subsubsection{Protocolli e sicurezza}

La 015 non impone un protocollo, ma indica i due più diffusi tra quelli
basati su YANG: NETCONF~\cite{rfc6241}, tipicamente tra apparato e
controller, transazionale e trasportato su SSH o TLS con codifica XML; e
RESTCONF~\cite{rfc8040}, tipicamente sulla NBI dei controller, basato su
HTTP con operazioni CRUD e codifica JSON o XML.

Sul piano della sicurezza (clausola 6) la specifica distingue due casi
all'installazione di un nuovo sistema. Se il sistema entra in un'area
sicura già presente nella rete, il canale verso il controller può essere
protetto con chiavi QKD prodotte da link esistenti. Se invece si aggiunge
un nuovo luogo sicuro, serve uno schema iniziale (un'identità fidata o un
segreto pre-condiviso, usato solo durante l'installazione) per presentare
il nuovo SDN agent al controller; dopo la prima autenticazione si passa a
chiavi QKD. È lo stesso principio di key growing visto in apertura.

% ---------------------------------------------------------------------
\subsubsection{Dal dominio singolo all'orchestrazione}

La 015 risolve il controllo di \emph{un} dominio QKD da parte di \emph{un}
controller. Restano fuori tre problemi che la GS~QKD~002 aveva già fatto
intravedere. Primo, un operatore può avere più domini QKD, magari di vendor
diversi, ciascuno con il proprio controller. Secondo, le applicazioni che
consumano chiave possono stare in un'altra rete, per esempio nei nodi della
rete di trasporto ottico, controllata da un controller diverso. Terzo,
l'operatore vuole gestire la QKD con gli stessi strumenti di orchestrazione
che usa per il resto della rete. Serve quindi un livello superiore e
un'interfaccia standard per raggiungerlo: è il compito della GS~QKD~018.
